\documentclass{article}
\usepackage[utf8]{inputenc}
\usepackage{amsmath}
\usepackage{geometry}
\geometry{margin=2cm}
\begin{document}

\section*{Questão 159 — ENEM 2024 — Matemática}

\subsection*{Enunciado}
Uma tubulação despeja sempre o mesmo volume de água por unidade de tempo em uma caixa-d’água, o que significa dizer que a vazão de água nessa tubulação é constante. Na junção dessa tubulação com a caixa-d’água, está instalada uma membrana de filtragem cujo objetivo é filtrar eventuais impurezas presentes na água, combinado a um bom fluxo de água. O fluxo ($\varphi$) de água através da superfície da membrana é diretamente proporcional à vazão de água na tubulação, medida em mililitro por segundo, e inversamente proporcional à área da superfície da membrana, medida em centímetro quadrado.

A unidade de medida adequada para descrever o fluxo ($\varphi$) de água que atravessa a superfície da membrana é:

\begin{figure}[ht]
\centering
\begin{tabular}{cc}
\textbf{A} & $\mathrm{mL \cdot s \cdot cm^{2}}$ \\
\includegraphics[height=1.5cm]{imagem_alternativa_A.png} & 
\textbf{B} \\[4pt]
& $\mathrm{\frac{mL}{s} \cdot cm^{2}}$ \\
\includegraphics[height=1.5cm]{imagem_alternativa_B.png} & 
\textbf{C} \\[4pt]
& $\mathrm{\frac{mL}{cm^{2} \cdot s}}$ \\
\includegraphics[height=1.5cm]{imagem_alternativa_C.png} & 
\textbf{D} \\[4pt]
& $\mathrm{\frac{cm^{2} \cdot s}{mL}}$
\includegraphics[height=1.5cm]{imagem_alternativa_D.png}
\end{tabular}
\end{figure}

\section{Solução Técnica}

A questão apresenta uma tubulação com vazão constante de água, expressa em mililitros por segundo (mL/s), e uma membrana com área em centímetro quadrado (cm$^2$). O fluxo $\varphi$ da água através da membrana é diretamente proporcional à vazão e inversamente proporcional à área da membrana. Assim, temos:

\begin{equation}
\varphi \propto \frac{\text{vazão}}{\text{área}}.
\end{equation}

Este raciocínio indica que a unidade do fluxo será dada pela unidade da vazão dividida pela unidade da área:

\begin{equation}
[\varphi] = \frac{\mathrm{mL/s}}{\mathrm{cm}^2} = \mathrm{\frac{mL}{cm^2 \cdot s}}.
\end{equation}

Comparando com as alternativas, a única que apresenta essa unidade é a alternativa \textbf{C}.

\subsection*{Passos da Resolução}
\begin{enumerate}
  \item Identificar a unidade da vazão: mililitros por segundo (mL/s).
  \item Compreender que o fluxo é proporcional à vazão e inversamente proporcional à área da membrana (cm$^2$).
  \item Escrever a expressão do fluxo: $\varphi = \frac{\text{vazão}}{\text{área}}$.
  \item Determinar a unidade do fluxo dividindo as unidades: $\frac{\mathrm{mL/s}}{\mathrm{cm}^2} = \mathrm{mL/(cm^2 \cdot s)}$.
  \item Escolher a alternativa que apresenta essa unidade, que é a \textbf{C}.
\end{enumerate}

\section*{Fórmulas Utilizadas}
\begin{equation}
\varphi \propto \frac{\text{vazão}}{\text{área}}.
\end{equation}

\section*{Evidência Visual Utilizada}
A análise considerou as imagens das unidades nas alternativas, especialmente a alternativa \textbf{C} que apresenta claramente a unidade $\mathrm{mL/(cm^2 \cdot s)}$.

\section{Explicação Didática}

\subsection*{Explicação Resumida}
O fluxo de água $\varphi$ é calculado dividindo-se a vazão da água (volume por tempo) pela área da membrana, pois $\varphi$ é diretamente proporcional à vazão e inversamente proporcional à área. Assim, a unidade do fluxo é mililitro por segundo por centímetro quadrado, ou seja, $\mathrm{mL/(cm^{2} \cdot s)}$.

\subsection*{Passo a Passo}
\begin{enumerate}
  \item Identificar que a vazão está em mililitros por segundo (mL/s).
  \item Reconhecer que o fluxo $\varphi$ é proporcional à vazão e inversamente proporcional à área da membrana (cm$^2$).
  \item Expressar matematicamente: $\varphi = \frac{\text{vazão}}{\text{área}}$.
  \item Determinar a unidade: dividir as unidades da vazão pela área, obtendo $\mathrm{mL/(cm^{2} \cdot s)}$.
  \item Comparar as unidades com as alternativas para identificar a correta.
\end{enumerate}

\subsection*{Erros Comuns}
\begin{itemize}
  \item Confundir multiplicação com divisão das unidades, o que leva a unidades incorretas.
  \item Inverter a proporcionalidade, multiplicando pela área em vez de dividir.
  \item Ignorar a combinação correta das unidades de volume, tempo e área.
  \item Não observar a ordem e os sinais das unidades na expressão final.
\end{itemize}

\subsection*{Dicas de Ensino}
\begin{itemize}
  \item Demonstre visualmente a diferença entre grandezas diretamente e inversamente proporcionais.
  \item Proponha exercícios de análise dimensional passo a passo.
  \item Incentive a verificação da consistência das unidades para evitar erros.
  \item Use exemplos do cotidiano para ilustrar as relações de proporção entre vazão, área e fluxo.
  \item Peça aos alunos para escreverem as relações em palavras antes de passarem para a matemática.
\end{itemize}

\section*{Erros Comuns na Resolução}
\begin{itemize}
  \item Confundir multiplicação com divisão das unidades, levando a resultados errados.
  \item Inverter a proporcionalidade entre fluxo e área, tratando a relação inversa como direta.
  \item Esquecer de combinar corretamente as unidades de volume (mL), tempo (s) e área (cm$^2$).
  \item Desconsiderar a ordem correta das unidades na expressão final do fluxo.
\end{itemize}

\section*{Dicas para o Ensino}
\begin{itemize}
  \item Utilizar exemplos visuais para mostrar grandezas diretamente e inversamente proporcionais.
  \item Desenvolver a habilidade de análise dimensional com atividades práticas.
  \item Estimular a verificação contínua da consistência das unidades durante a resolução.
  \item Relacionar conceitos com situações do dia a dia para melhor compreensão.
  \item Incentivar a formulação verbal das relações matemáticas antes de trabalhar com unidades.
\end{itemize}

\section{Análise dos Distratores}

\begin{itemize}
  \item \textbf{Alternativa A}: A unidade apresentada é $\mathrm{mL \cdot s \cdot cm^{2}}$. Isso indica uma confusão entre multiplicação e divisão das unidades, o que é um erro conceitual em análise dimensional. Embora as unidades envolvidas estejam relacionadas ao problema, a forma como estão combinadas está incorreta.

  \item \textbf{Alternativa B}: A unidade dada é $\mathrm{\frac{mL}{s} \cdot cm^{2}}$. Aqui a área foi multiplicada pela vazão, invertendo a relação inversa correta entre fluxo e área, o que constitui um erro conceitual na proporcionalidade.

  \item \textbf{Alternativa D}: Apresenta a unidade $\mathrm{\frac{cm^{2} \cdot s}{mL}}$, invertendo as grandezas no numerador e denominador sem lógica física. Este é um erro na formulação da expressão do fluxo e na análise dimensional.
\end{itemize}

Apesar de plausíveis a partir da relação com as grandezas dadas, todas essas alternativas erradas falham em respeitar a proporcionalidade correta entre vazão, área e fluxo.

\section{Perfil do Item}

\begin{itemize}
  \item \textbf{Habilidade Primária}: Análise dimensional.
  \item \textbf{Habilidades Secundárias}: Razão e proporção; compreensão de grandezas físicas.
  \item \textbf{Nível Cognitivo}: Médio.
  \item \textbf{Dificuldade Estimada}: Média.
  \item \textbf{Requisitos}: Interpretação visual das unidades; manipulação algébrica simples.
  \item \textbf{Notas}: Envolve entendimento de grandezas diretamente e inversamente proporcionais e exige atenção às unidades físicas e sua correta combinação.
\end{itemize}


\end{document}